\scanpage{38}
% archive sheet number: 37
% slanted
\margin{\handwritten{\unsure{Nouveau|Nouveaux} Chapitre}}
\handwritten{1}) Programme de la théorie de Dieudonné sur une base \inserted{$S$} où $p$ est localement nilpotent.

a) \uline{Foncteur de Dieudonné}. On \deleted{va} définiraun foncteur
\margin{\handwritten{$\mathbb{D}^*$}}
\[
\underline{D}^{*}:\quad \mathrm{BT}(S)^o\handwritten{\longrightarrow}\mathrm{Crisloclib}(S)\quad,
\]
compatible aux images inverses (donc foncteur cartésien sur la catégorie fibrée des groupes de BT sur des bases où $p$ est loc\deleted{\unsure{,|;}}\handwritten{\unsure{.|;}} nilp.). Si $S_{\unsure{o|0}}=\mathrm{Var}(p)$ on a une équivalence de catégories (cf.\qquad plus haut)
\[
\mathrm{Crisloclib}(S)\handwritten{\approxeq}\mathrm{Crisloclib}(S_o)
\]
de sorte que la donnée de $\underline{D}^{*}$ équivaut à la donnée de
\[
\underline{D}^{*}:\quad\mathrm{BT}(S_o)^o\handwritten{\longrightarrow}\mathrm{Crisloclib}(S)_o\quad,
\]
i.e. on est ramené \deleted{à la car. $p$} au cas $S$ de car. $p$. \deleted{Pour connaitre, da}

Dans le cas général, $\underline{D}^{*}(G)$ ne dépend que de \deleted{$\unsure{\underline{D}^*|D^*}$} $G_o=G\mathbin{\mathrm{x}}_S S_o$.

\deleted{Utilis} Plaçons nous sur $\deleted{\unsure{X}}S_o$ de car. $p$. Utilisant $\underline{f}_G$ et $\underline{v}_G$ et la compatibilité de $\underline{D}^{*}$ aux images inverses, on voit que \deleted{le foncteur} $\underline{D}^{*}(G_o)$ \deleted{peut}\inserted{$=M$} est naturellement muni de morphismes $F_M,V_M$ \deleted{satisfaisant}
\diagram{\begin{tikzcd}[ampersand replacement=\&]
M^{(p)}=\underline{f}_S^{*}(M) \arrow[r,shift left,"\handwritten{F_M}"] \& M \arrow[l,shift left,"\handwritten{V_M}"]
\end{tikzcd}}
satisfaisant
\[
V_M F_M=p\operatorname{id}_M,\quad F_M V_M=p\operatorname{id}_{\deleted{\unsure{m}}M^{(p)}}\quad.
\]
\deleted{\unsure{On}}On appelle un tel triple $(M,F_M,V_M)$ un (F-V-\uline{cristal}\inserted{\handwritten{ ou \uline{cristal de Dieudonné})}}, et on a donc un foncteur
\[
\underline{D}^{*}:\quad\mathrm{BT}(S_o)\handwritten{\longrightarrow}\text{F-V-Cris}(S_o)\quad .
\]
C'est ce foncteur qui mérite le nom de \uline{foncteur de Dieudonné} . On \deleted{\unsure{pura|poura} prouver que} voit que sa formation commute à tout changement de\handwritten{\inserted{ }}base. De plus, on prouve que pour $S_o$ spectre d'un corps parfait, il coïncide \deleted{avec} à isomorphisme canonique près avec le foncteur de Dieudonné habituel, qui est alors une équivalence de catégories, de sorte que dans ce cas la
